\section{Method}
\label{sec:method}

We consider on-policy reinforcement learning with a stochastic flow policy
$\pi_\theta$.  At each environment step, the agent observes state
$\mathbf{s}\in\R^{d_s}$ and produces action
$\mathbf{a}\in[-1,1]^{d_a}$ by running a $K$-step stochastic generative
chain.

% ------------------------------------------------------------------
\subsection{Stochastic Flow Policy}
\label{sec:method:flow}

The generative chain proceeds from pure noise to an action as follows.
Let $\Delta t = 1/K$ and let $\sigma > 0$ be the per-step noise scale.

\begin{enumerate}[leftmargin=*,nosep]
  \item Sample initial noise: $\mathbf{z}_K \sim \mathcal{N}(\mathbf{0},\mathbf{I})$.
  \item For $k = K, K{-}1, \ldots, 1$, compute the next latent:
\end{enumerate}
\begin{align}
  \mathbf{z}_{k-1}
    &= \mathbf{z}_k
       + \Delta t \cdot v_\theta(\mathbf{s}, \mathbf{z}_k, k)
       + \sigma\,\boldsymbol{\epsilon}_k,
  \qquad \boldsymbol{\epsilon}_k \sim \mathcal{N}(\mathbf{0},\mathbf{I}).
  \label{eq:flow_step}
\end{align}
\begin{enumerate}[leftmargin=*,nosep,resume]
  \item Apply a squashing function to obtain the action:
        $\mathbf{a} = \tanh(\mathbf{z}_0)$.
\end{enumerate}

Each transition has a tractable Gaussian conditional:
\begin{align}
  p_\theta(\mathbf{z}_{k-1} \mid \mathbf{z}_k, \mathbf{s}, k)
    &= \mathcal{N}\!\bigl(
         \mathbf{z}_{k-1};\;
         \mathbf{z}_k + \Delta t\, v_\theta(\mathbf{s}, \mathbf{z}_k, k),\;
         \sigma^2 \mathbf{I}
       \bigr).
  \label{eq:per_step_density}
\end{align}

The holistic log-density of the final action is:
\begin{align}
  \log p_\theta(\mathbf{a}\mid\mathbf{s})
    &= \sum_{k=1}^{K}
       \log p_\theta(\mathbf{z}_{k-1}\mid\mathbf{z}_k,\mathbf{s},k)
       + \log\!\left|\det\frac{\partial\,\tanh(\mathbf{z}_0)}
                               {\partial\,\mathbf{z}_0}\right|.
  \label{eq:holistic_logprob}
\end{align}

% ------------------------------------------------------------------
\subsection{Holistic PPO (Baseline)}
\label{sec:method:holistic}

Standard PPO computes a single importance-sampling ratio per environment
transition $(s_t, a_t)$:
\begin{align}
  r_t
    &= \frac{p_\theta(\mathbf{a}_t\mid\mathbf{s}_t)}
            {p_{\theta_{\text{old}}}(\mathbf{a}_t\mid\mathbf{s}_t)},
  \label{eq:holistic_ratio}
\end{align}
and optimizes the clipped surrogate:
\begin{align}
  L^{\text{CLIP}}(\theta)
    &= \E_t\Bigl[
         \min\!\bigl(
           r_t\, \hat{A}_t,\;
           \clip(r_t, 1{-}\epsilon, 1{+}\epsilon)\,\hat{A}_t
         \bigr)
       \Bigr],
  \label{eq:ppo_clip}
\end{align}
where $\hat{A}_t$ is the GAE advantage estimate and $\epsilon$ is the clip
parameter (typically $0.2$).

Because $r_t$ is the product of $K$ per-step ratios (plus a correction),
it can be very large or very small, making clipping less effective for
large $K$.

% ------------------------------------------------------------------
\subsection{Per-Step PPO}
\label{sec:method:perstep}

We propose to apply PPO's clipped surrogate \emph{independently at each
generative step}.  Define the per-step ratio:
\begin{align}
  r_{t,k}
    &= \frac{p_\theta(\mathbf{z}_{t,k-1}\mid\mathbf{z}_{t,k},\mathbf{s}_t,k)}
            {p_{\theta_{\text{old}}}(\mathbf{z}_{t,k-1}\mid\mathbf{z}_{t,k},\mathbf{s}_t,k)}.
  \label{eq:perstep_ratio}
\end{align}

Each step $k$ is assigned an advantage $A_{t,k} = w_k \cdot \hat{A}_t$, where
$w_k$ is a credit-assignment weight satisfying $\sum_k w_k = 1$.  The per-step
clipped surrogate is:
\begin{align}
  L^{\text{PS-CLIP}}(\theta)
    &= \E_t\!\left[
         \sum_{k=1}^{K}
         \min\!\bigl(
           r_{t,k}\, A_{t,k},\;
           \clip(r_{t,k}, 1{-}\epsilon, 1{+}\epsilon)\, A_{t,k}
         \bigr)
       \right].
  \label{eq:perstep_clip}
\end{align}

By clipping each $r_{t,k}$ individually, we prevent any single step from
producing an excessively large update, even when the holistic ratio $r_t$
would be far from~$1$.

% ------------------------------------------------------------------
\subsection{Weighted Credit Assignment}
\label{sec:method:weights}

We study three weighting schemes for the per-step advantages
$A_{t,k} = w_k \cdot \hat{A}_t$:

\paragraph{Uniform.}
All steps receive equal credit:
\begin{align}
  w_k &= \frac{1}{K}, \qquad \forall\, k.
  \label{eq:uniform}
\end{align}

\paragraph{Learned Global.}
A single set of learnable logits $\boldsymbol{\alpha}\in\R^K$ produces
step-dependent but state-independent weights:
\begin{align}
  w_k &= \softmax(\boldsymbol{\alpha})_k.
  \label{eq:learned_global}
\end{align}
The logits $\boldsymbol{\alpha}$ are trained end-to-end alongside $\theta$.

\paragraph{State-Dependent.}
A small network $h_\phi$ takes the current state and latent as input and
produces per-step weights that depend on the rollout:
\begin{align}
  w_{t,k} &= \softmax\!\bigl(
                h_\phi(\mathbf{s}_t, \mathbf{z}_{t,k}, k)
              \bigr)_k.
  \label{eq:state_dep}
\end{align}
Here $\phi$ denotes the parameters of the weighting network, which are
jointly optimized with $\theta$.

% ------------------------------------------------------------------
\subsection{Intra-Chain Temporal Credit Assignment}
\label{sec:method:intra}

The weighting schemes above split the environment advantage $\hat{A}_t$
into per-step shares using fixed or learned proportions, but these shares
are heuristic---they do not reflect which steps \emph{actually improved}
the latent toward a good action.

We propose \textbf{intra-chain temporal credit assignment}: we learn a
value function $V_\psi(\mathbf{s}, \mathbf{z}_k, k)$ that estimates the
expected action quality at each internal step, and compute genuine
per-step advantages via temporal differencing within the chain:
\begin{align}
  A_{t,k}^{\text{intra}}
    &= V_\psi(\mathbf{s}_t, \mathbf{z}_{t,k-1}, k{-}1)
     - V_\psi(\mathbf{s}_t, \mathbf{z}_{t,k}, k).
  \label{eq:intra_adv}
\end{align}
Intuitively, $A_{t,k}^{\text{intra}}$ measures how much step $k$
improved the latent---a step that moves the latent toward a
high-value region of the chain receives positive credit, while a step
that degrades quality receives negative credit.

To ground these intra-chain advantages in the actual reward signal, we
normalize them to sum to the environment advantage:
\begin{align}
  \tilde{A}_{t,k}
    &= A_{t,k}^{\text{intra}} \cdot
       \frac{\hat{A}_t}
            {\sum_{j=1}^{K} A_{t,j}^{\text{intra}} + \delta},
  \label{eq:intra_norm}
\end{align}
where $\delta$ is a small constant for numerical stability.

The intra-chain value function is trained with the loss:
\begin{align}
  L^{\text{intra}}(\psi)
    &= \E_t\!\left[\bigl(
         V_\psi(\mathbf{s}_t, \mathbf{z}_{t,0}, 0)
       - V_\psi(\mathbf{s}_t, \mathbf{z}_{t,K}, K)
       - \hat{A}_t
       \bigr)^2\right],
  \label{eq:intra_loss}
\end{align}
which ensures the total value change across the chain predicts the
environment advantage.  The per-step PPO loss (\Cref{eq:perstep_clip})
then uses $\tilde{A}_{t,k}$ in place of $w_k \cdot \hat{A}_t$.

This approach extends temporal credit assignment \emph{inside} the
generative chain, analogous to how GAE assigns credit across
environment timesteps.  Unlike the weight-splitting schemes
(\Cref{sec:method:weights}), intra-chain credit is state- and
latent-dependent, grounded in learned value estimates, and need not
sum to a fixed proportion.

% ------------------------------------------------------------------
\subsection{Algorithm Summary}
\label{sec:method:algorithm}

\Cref{alg:perstep_ppo} summarizes the per-step PPO training loop.

\begin{algorithm}[t]
\caption{Per-Step PPO for Stochastic Flow Policies}
\label{alg:perstep_ppo}
\KwIn{Initial policy $\theta_0$, weight parameters (depending on scheme),
      clip parameter $\epsilon$, number of flow steps $K$, noise scale $\sigma$}
\For{iteration $= 1, 2, \ldots$}{
  Collect rollout batch using $\pi_{\theta_{\text{old}}}$\;
  \For{each transition $(\mathbf{s}_t, \mathbf{a}_t, \hat{A}_t)$}{
    Store intermediate latents $\{\mathbf{z}_{t,k}\}_{k=0}^{K}$\;
  }
  Compute GAE advantages $\{\hat{A}_t\}$\;
  \For{PPO epoch $= 1, \ldots, N_{\text{epochs}}$}{
    \For{minibatch of transitions}{
      \For{$k = K, \ldots, 1$}{
        Compute $r_{t,k}$ via \Cref{eq:perstep_ratio}\;
        Compute $w_k$ (or $w_{t,k}$) via chosen scheme\;
        $A_{t,k} \leftarrow w_k \cdot \hat{A}_t$\;
      }
      $L \leftarrow L^{\text{PS-CLIP}}$ \hfill (\Cref{eq:perstep_clip})\;
      Update $\theta$ (and $\phi$ or $\boldsymbol{\alpha}$ if applicable)
        via gradient ascent on $L$\;
    }
  }
}
\end{algorithm}
